\documentclass[11pt]{article}

% --- Packages for math & formatting ---
\usepackage{amsmath, amssymb, amsthm} % advanced math
\usepackage{listings}                  % code formatting
\usepackage{xcolor}                     % color for code
\usepackage{graphicx}                   % include figures
\usepackage{hyperref}                   % clickable links
\usepackage{authblk}                    % author affiliations
\usepackage[margin=1in]{geometry}
\usepackage{fvextra} % better verbatim with line wrapping
\usepackage{float}

% --- Configure code display ---
\lstset{
  language=R,
  basicstyle=\ttfamily\small,
  numbers=left,
  stepnumber=1,
  numbersep=5pt,
  frame=single,
  backgroundcolor=\color{gray!10},
  breaklines=true,
  captionpos=b
}

% --- Theorem environment ---
\newtheorem{theorem}{Theorem}
\newtheorem{lemma}{Lemma}

% --- New commands ---
\newcommand*\mean[1]{\overline{#1}}
\newcommand{\var}{\text{var}}
\newcommand{\cov}{\text{cov}}
\newcommand{\cor}{\text{cor}}
\newcommand{\Rp}{\text{Re}}
\newcommand{\E}{\text{E}}
\newcommand{\ltsgr}{\text{ltsgr}}
\newcommand{\expit}{\text{expit}}
\newcommand{\logit}{\text{logit}}

% --- Header material ---
\title{How to fit models with species occurrence data using \texttt{xsdm}}
\author[1]{Daniel C. Reuman}
\author[1]{Angel Lu\'{i}s Robles Fern\'{a}ndez}
\affil[1]{Department of Ecology and Evolutionary Biology and Center for Ecological Research,
University of Kansas}
\date{}

% --- Make Sweave code/output blocks obey text margins and wrap nicely ---
\DefineVerbatimEnvironment{Sinput}{Verbatim}{
  fontsize=\small,
  xleftmargin=0pt,
  xrightmargin=0pt,
  breaklines=true,
  breakanywhere=true,
  commandchars=\\\{\}
}

\DefineVerbatimEnvironment{Soutput}{Verbatim}{
  fontsize=\small,
  xleftmargin=0pt,
  xrightmargin=0pt,
  breaklines=true,
  breakanywhere=true,
  commandchars=\\\{\}
}

% Optional: if you have Scode blocks too
\DefineVerbatimEnvironment{Scode}{Verbatim}{
  fontsize=\small,
  xleftmargin=0pt,
  xrightmargin=0pt,
  breaklines=true,
  breakanywhere=true,
  commandchars=\\\{\}
}

% --- start the doc ---
\usepackage{Sweave}
\begin{document}
\Sconcordance{concordance:howToFit.tex:howToFit.Rnw:1 79 1 1 0 2 1 1 7 9 0 1 2 27 1 1 %
2 1 0 4 1 7 0 1 1 5 0 2 1 5 0 1 1 5 0 1 1 3 0 1 2 2 1 1 2 1 0 1 1 7 0 1 %
1 5 0 2 1 5 0 1 1 6 0 1 2 4 1 1 2 1 0 2 1 5 0 1 1 5 0 2 1 5 0 1 1 6 0 1 %
2 1 1 1 3 1 0 1 1 5 0 1 1 5 0 1 1 11 0 1 1 6 0 1 2 4 1 1 2 1 0 4 1 3 0 %
1 2 3 1 1 2 1 0 4 1 3 0 1 2 1 1 1 2 1 0 7 1 3 0 1 2 6 1 1 2 1 0 2 1 5 0 %
1 1 6 0 1 2 1 1 1 2 1 0 1 1 6 0 1 2 27 1 1 2 1 0 5 1 1 8 13 0 1 2 1 1 1 %
2 7 0 1 2 54 1 1 4 3 0 1 1 8 0 1 2 2 1 23 0 1 2 6 0 1 2 3 1 1 2 7 0 1 2 %
13 1 1 8 36 0 1 2 5 1 1 2 1 0 2 1 18 0 1 2 5 1 1 2 1 0 1 10 12 0 1 2 7 %
1 1 2 1 0 2 1 7 0 3 1 7 0 1 2 9 1 1 2 1 0 2 1 25 0 1 2 25 1 1 2 1 0 1 8 %
7 0 1 1 7 0 1 2 15 1 1 3 2 0 3 1 1 17 19 0 1 2 1 1 1 3 2 0 1 1 7 0 4 1 %
6 0 1 3 1 0 1 1 1 6 5 0 1 1 7 0 1 2 3 1 1 3 2 0 1 1 7 0 4 1 6 0 1 3 1 0 %
1 1 1 6 5 0 1 1 7 0 1 2 4 1 1 4 3 0 1 1 5 0 1 3 1 0 1 1 5 0 2 1 6 0 1 2 %
24 1 1 13 12 0 1 1 5 0 1 1 12 0 1 2 5 1 1 6 5 0 1 2 4 0 1 2 12 1 1 2 19 %
0 1 1 3 0 1 2 4 1 1 2 7 0 1 2 14 1 1 2 1 0 1 1 1 15 14 0 1 1 3 0 1 2 4 %
1 1 9 8 0 1 1 1 3 2 0 1 1 24 0 1 2 2 1 1 2 1 0 1 9 7 0 1 7 5 0 1 7 5 0 %
1 7 5 0 1 7 5 0 1 7 5 0 1 7 5 0 1 7 5 0 1 7 5 0 1 5 3 0 1 5 7 0 1 2 14 %
1 1 3 2 0 2 1 1 3 1 0 1 2 1 0 1 4 2 0 1 4 2 0 1 2 1 0 1 5 3 0 1 3 1 0 3 %
1 5 0 1 1 5 0 1 1 6 0 1 2 8 1 1 3 2 0 2 1 1 3 1 0 3 1 1 2 1 0 1 4 2 0 1 %
4 2 0 1 2 1 0 1 5 3 0 1 3 1 0 3 1 5 0 1 1 5 0 1 1 6 0 1 2 17 1 1 2 1 0 %
1 2 1 5 7 0 1 2 14 1 1 2 1 0 1 5 8 0 1 3 5 1}


\begin{Schunk}
\begin{Sinput}
> knitr::opts_chunk$set(
+   echo=TRUE,
+   message=FALSE,
+   warning=FALSE,
+   cache=FALSE
+ )
\end{Sinput}
\end{Schunk}

\maketitle

\begin{abstract}
\noindent After reading the document entitled "The xsdm model", this document
introduces the statistically and computationally sophisticated user to the 
main functions of the \texttt{xsdm} package. The package implements a frequentist 
analysis of the xsdm model. This document should get the user to the point
of maximizing the likelihood of the model and profiling in ``easy'' cases, 
with callouts to
several troubleshooting documents about what to do when the process does
not go smoothly. Alternative fitted models can also be compared by AIC or another
criterion. The example worked out here is for a virtual species (i.e., generated 
data), so we also demonstrate that fitting an xsdm model can recover the true
parameters if model mis-specification is not an issue.
\end{abstract}

\section{Introduction to the running example used throughout this document}

This manual uses a small built-in example included with the \texttt{xsdm} R 
package to illustrate a baseline workflow. The example contains (i) an
occurrence table for a virtual species, and (ii) two environmental
\texttt{terra} raster time series representing CHELSA-derived 
bioclimatic variables for a region of interest over 39 years.

Load the data and get oriented to the first bioclimatic variable, which is mean annual temperature
in units of $100 \times ^\circ$C for 39 years for a particular region of interest. 
And make it units of $^\circ$C: 
\begin{Schunk}
\begin{Sinput}
> library(xsdm)
> library(terra)